\exercisesection{4}

\begin{enumerate}
\def\labelenumi{\arabic{enumi}.}
\tightlist
\item
  设 A 为整闭 Noether 整环，V 为 A 的分式域上的有限秩向量空间。证明：若 \((\mathbf{M}_{\lambda})\) 是 V 中都包含同一个格 N 的自反格的任一族，则格 \(M = \bigcap_{\lambda} M_{\lambda}\) 是自反的（考虑对偶格 \(M_{\lambda}^{*}\)）。
\end{enumerate}

* 2. 设 A 为整闭 Noether 整环，E、F 为两个有限生成 A-模。假设 E 是无扭的，\(\mathrm{Hom}_{\mathrm{A}}(\mathrm{E},\mathrm{F})\) 是自反的，且 \(\mathrm{Ext}_{\mathrm{A}}^{1}(\mathrm{E},\mathrm{F}) = 0\)。证明 F 是自反的。（先证明 F 是无扭的；若 \(\mathrm{T} = \mathrm{F}^{**}/c_{\mathrm{F}}(\mathrm{F})\)，利用正合列以两种方式计算 \(\mathrm{Ass}(\mathrm{Hom}_{\mathrm{A}}(\mathrm{E},\mathrm{T}))\)：

\[
0 \rightarrow \operatorname{Hom} (\mathrm{E}, \mathrm{F}) \rightarrow \operatorname{Hom} (\mathrm{E}, \mathrm{F} ^ {* *}) \rightarrow \operatorname{Hom} (\mathrm{E}, \mathrm{T}) \rightarrow 0
\]

以及第 IV 章 § 1 第 4 节的命题 10。） *

\begin{enumerate}
\def\labelenumi{\arabic{enumi}.}
\setcounter{enumi}{2}
\item
  设 A 为整闭 Noether 整环，K 为其分式域，M 为 K 上某有限秩向量空间的自反格，L 为包含 M 的 V 的自由格。证明存在 V 的自由格 \(L_{1}\) 使 \(M = L \cap L_{1}\)。（考虑使 \(L_{p} \neq M_{p}\) 的 A 的高度为 1 的素理想 p 所成的有限集 I，以及主理想整环 \(S^{-1}A\)，其中 \(S = \bigcap_{p \in I} (A - p)\)；证明存在 V 的自由格 \(L_{0}\)，使 \(M_{p} = (L_{0})_{p}\) 对所有 \(p \in I\) 成立，且存在 \(s \in S\)，使 \(M \subset s^{-1}L_{0} \subset L_{1}\)。）
\item
  设 k 为域，A 为多项式环 \(k[X, Y]\)。
\end{enumerate}

\begin{enumerate}
\def\labelenumi{(\alph{enumi})}
\item
  设 \((e_1, e_2)\) 为 A-模 \(A^2\) 的典范基，E 为 \(A^2\) 的由 \((X - Y)e_1\)、\(e_1 + Xe_2\) 与 \(e_1 + Ye_2\) 生成的 A-子模；设 F 为 E 的由 \((X - Y)^2e_1\) 生成的单生成子模。证明 A-模 \(M = E / F\) 不是其扭子模与一个无扭模的直和。
\item
  证明扭 A-模 A/AXY 不是形如 \(A/P_{i}^{n_{i}}\) 的单生成子模的直和，其中诸 \(P_{i}\) 是高度为 1 的素理想。
\end{enumerate}

\begin{enumerate}
\def\labelenumi{\arabic{enumi}.}
\setcounter{enumi}{4}
\tightlist
\item
  设 A 为整闭 Noether 整环，M 为有限生成 A-模。证明：若 a、b 是 A 的两个理想，则 A-模
\end{enumerate}

\[
((a \mathbf {M}) \cap (\mathfrak {b} \mathbf {M})) / (a \cap \mathfrak {b}) \mathbf {M}
\]

是伪零的；给出一个使它不为零的例子。

\begin{enumerate}
\def\labelenumi{\arabic{enumi}.}
\setcounter{enumi}{5}
\tightlist
\item
  设 \(k\) 为域，\(\mathbf{B}\) 为多项式环 \(k[\mathbf{X},\mathbf{Y}]\)，\(\mathbf{A}\) 为子模 \(k[\mathbf{X}^2,\mathbf{XY},\mathbf{Y}^2]\)（包含于 \(\mathbf{B}\)）。
\end{enumerate}

\begin{enumerate}
\def\labelenumi{(\alph{enumi})}
\item
  证明 A 是整闭 Noether 整环，且 B 是有限生成 A-模。
\item
  证明 A 的理想 \(p = AX^{2} + AXY\) 是高度为 1 的素理想（从而是除性的），但 B-模 \(p \otimes_{A} B\) 是非零 \(\neq 0\) 的扭模，从而不是自反的；B 的理想 pB 不是除性的，典范映射 \(p \otimes_{A} B \to pB\) 不是单射，且 p 不是平坦 A-模。
\end{enumerate}

\begin{enumerate}
\def\labelenumi{\arabic{enumi}.}
\setcounter{enumi}{6}
\tightlist
\item
  设 A 为整闭 Noether 整环，E 为有限生成无扭 A-模，\(\mathbf{E}^*\) 为其对偶模。
\end{enumerate}

\begin{enumerate}
\def\labelenumi{(\alph{enumi})}
\item
  证明典范同态 \(\mathbf{E}^{*}\otimes_{\mathbf{A}}\mathbf{E}\to \mathrm{End}_{\mathbf{A}}(\mathbf{E})\) 是伪同构。
\item
  由 (a) 推出：E 为投射 A-模的充分必要条件是 \(E^{*} \otimes_{A} E\) 为自反 A-模。（注意，若 \(E^{*} \otimes_{A} E\) 自反，则典范同态 \(E^{*} \otimes_{A} E \to \operatorname{End}_{A}(E)\) 是双射。）
\end{enumerate}

¶* 8. 设 A 为整闭 Noether 整环，\(\mathbf{M}_1, \mathbf{M}_2\) 为两个有限生成 A-模。

\begin{enumerate}
\def\labelenumi{(\alph{enumi})}
\item
  证明 A-模 \(\operatorname{Tor}_i^{\mathbf{A}}(\mathbf{M}_1, \mathbf{M}_2)\)、\(\operatorname{Ext}_{\mathbf{A}}^{i}(\mathbf{M}_1, \mathbf{M}_2)\) 对 \(i \geqslant 2\) 都是伪零的（归结为 \(\mathbf{A}\) 是主理想整环的情形）。
\item
  若 \(\mathbf{M}_1\) 是无扭的，证明 \(\operatorname{Tor}_1^{\mathrm{A}}(\mathbf{M}_1, \mathbf{M}_2)\) 与 \(\operatorname{Ext}_A^1(\mathbf{M}_1, \mathbf{M}_2)\) 是伪零的（同样的方法）。
\item
  若 \(\mathbf{M}_1\) 是扭 A-模，证明（用第 5 节的记号）：
\end{enumerate}

\[
\chi (\mathbf {M} _ {1} \otimes_ {\mathbf {A}} \mathbf {M} _ {2}) - \chi (\operatorname{Tor} _ {1} ^ {\mathbf {A}} (\mathbf {M} _ {1}, \mathbf {M} _ {2})) = r (\mathbf {M} _ {2}) \chi (\mathbf {M} _ {1})
\]

\[
\chi (\operatorname{Hom} _ {\mathbf {A}} (\mathbf {M} _ {1}, \mathbf {M} _ {2})) - \chi (\operatorname{Ext} _ {\mathbf {A}} ^ {1} (\mathbf {M} _ {1}, \mathbf {M} _ {2})) = - r (\mathbf {M} _ {2}) \chi (\mathbf {M} _ {1})
\]

\[
\chi (\mathrm{Hom} _ {\mathrm{A}} (\mathrm{M} _ {2}, \mathrm{M} _ {1})) - \chi (\mathrm{Ext} _ {\mathrm{A}} ^ {1} (\mathrm{M} _ {2}, \mathrm{M} _ {1})) = r (\mathrm{M} _ {2}) \chi (\mathrm{M} _ {1}).
\]

\begin{enumerate}
\def\labelenumi{(\alph{enumi})}
\setcounter{enumi}{3}
\tightlist
\item
  证明对有限生成 A-模的每个有序对 \(M_{1}\)、\(M_{2}\)：
\end{enumerate}

\[
\begin{array}{r} c (\mathbf {M _ {1}} \otimes_ {\mathbf {A}} \mathbf {M _ {2}}) - c (\mathrm{Tor} _ {\mathbf {1}} ^ {\mathbf {A}} (\mathbf {M _ {1}}, \mathbf {M _ {2}})) = r (\mathbf {M _ {1}}) c (\mathbf {M _ {2}}) + r (\mathbf {M _ {2}}) c (\mathbf {M _ {1}}) \\ c (\mathrm{Hom} _ {\mathbf {A}} (\mathbf {M _ {1}}, \mathbf {M _ {2}})) - c (\mathrm{Ext} _ {\mathbf {A}} ^ {1} (\mathbf {M _ {1}}, \mathbf {M _ {2}})) = r (\mathbf {M _ {1}}) c (\mathbf {M _ {2}}) - r (\mathbf {M _ {2}}) c (\mathbf {M _ {1}}). \end{array}
\]

\begin{enumerate}
\def\labelenumi{\arabic{enumi}.}
\setcounter{enumi}{8}
\tightlist
\item
  设 k 为域，A 为多项式环 \(k[X,Y]\)。在 A-模 \(M = A^{2}\) 上考虑满足 \(f(e_{1}) = X\)、\(f(e_{2}) = Y((e_{1}, e_{2})\) 为典范基）的线性形式 f。证明 f 的核 L 是单生成自由 A-模，但商模 M/L（它同构于 A 的一个理想）不是自反的。
\end{enumerate}

¶ 10. 设 A 为交换环。对有限生成自由 A-模 \(L = A^n\) 的每个子模 R，用 \(c_1(R)\) 表示由诸 \(\langle x, x^* \rangle\) 生成的理想，其中 \(x\) 取遍 R，\(x^*\) 取遍对偶 \(L^*\)；记 \(c_k(R) = c_1\left(\operatorname{Im}\left(\bigwedge^k R\right)\right)\)，其中 \(\operatorname{Im}\left(\bigwedge^k R\right)\) 是 R 的 \(k\) 次外幂在 \(\bigwedge^k L\) 中的典范像。若 \(R_1 \subset R_2\) 是 L 的两个子模，则 \(c_k(R_1) \subset c_k(R_2)\) 对所有 \(k\) 成立。

\begin{enumerate}
\def\labelenumi{(\alph{enumi})}
\item
  设 M 为有限生成 A-模；M 同构于一个商模 L/R，其中对某个适当的 n 有 \(L = A^{n}\)。证明：理想序列 \((\mathfrak{a}_{k})_{k \geqslant 0}\)（满足 \(\mathfrak{a}_{k} = \mathfrak{c}_{n-k}(R)\)，这里 \(k \leqslant n\)；\(\mathfrak{a}_{k} = A\)，这里 k \textgreater{} n）与 M 表成 L/R 的方式无关。（先考虑 L/R 与 L/R' 同构的情形；然后注意，\(A^{n}/R\) 同构于 \(A^{n+h}/(R \times A^{h})\)，对所有 h \textgreater{} 0。）我们记 \(\mathfrak{d}_{k}(M) = \mathfrak{a}_{k}\)（对所有 \(k \geqslant 0\)），并称它们为与 M 相伴的行列式理想。于是 \(\mathfrak{d}_{k}(M) \subset \mathfrak{d}_{k+1}(M)\) 对 \(k \geqslant 0\) 成立。
\item
  若 \(\mathbf{A}\) 是主理想整环且 \(\mathbf{M} = \mathbf{L} / \mathbf{R}\)，其中 \(\mathbf{L}\) 是自由且有限生成的，证明理想 \(\mathfrak{c}_{k + 1}(\mathbb{R})(\mathfrak{c}_k(\mathbb{R}))^{-1}\) 是 \(\mathbf{R}\) 在 \(\mathbf{L}\) 中的不变因子。
\item
  设 \(r = \gamma(\mathbf{M})\) 为 \(\mathbf{M}\) 的生成系的最小基数，\(r_0\) 为最小整数 \(h\)，使 \(\mathfrak{d}_h(\mathbf{M}) = \mathbf{A}\)。证明 \(r_0 \leqslant r\)，并给出一个使 \(r_0 < r\) 的例子（取 \(\mathbf{A}\) 为 Dedekind 整环）。
\item
  若 \(a\) 是 \(\mathbf{M}\) 的零化子，证明 \(\mathfrak{d}_0(\mathbf{M}) \subset \mathfrak{a}\)，且 \(\mathfrak{a}^{r - k} \subset \mathfrak{d}_k(\mathbf{M})\) 对 \(k \leqslant r = \gamma (\mathbf{M})\) 成立。
\item
  设 \(\mathbf{M}\) 是子模 \(\mathbf{M}_i\)（\(1 \leqslant i \leqslant h\)）的直和。证明
\end{enumerate}

\[
\mathfrak {d} _ {k} (\mathbf {M}) = \sum \mathfrak {d} _ {k _ {1}} (\mathbf {M} _ {1}) \dots \mathfrak {d} _ {k _ {n}} (\mathbf {M} _ {n})
\]

其中和式取遍有限序列 \((k_{i})_{1\leqslant i\leqslant h}\)（满足 \(\sum_{i = 1}^{h}k_{i} = k.\)）

\begin{enumerate}
\def\labelenumi{(\alph{enumi})}
\setcounter{enumi}{5}
\tightlist
\item
  若 N 是 M 的有限生成子模，则 \(\mathfrak{d}_{k}(\mathbf{M}/\mathbf{N}) \subset \mathfrak{d}_{k}(\mathbf{M})\) 对所有 \(k \geqslant 0\) 成立。证明：
\end{enumerate}

\[
\sum_ {j = 0} ^ {k} \mathfrak {d} _ {j} (\mathrm{N}) \mathfrak {d} _ {k - j} (\mathrm{M} / \mathrm{N}) \subset \mathfrak {d} _ {k} (\mathrm{M}).
\]

\begin{enumerate}
\def\labelenumi{(\alph{enumi})}
\setcounter{enumi}{6}
\tightlist
\item
  设 K 为域，A 为多项式环 K{[}X, Y, Z{]}，m 为极大理想 AX + AY + AZ，q 为由 X、Y²、YZ 与 Z² 生成的理想。考虑 A-模 M = A/q 及其子模 N = m/q。证明 \(\mathfrak{d}_{0}(M) = \mathfrak{q}, \mathfrak{d}_{1}(M) = A, \mathfrak{d}_{0}(N) = m^{2}\) 且 \(\mathfrak{d}_{1}(N) = m\)。
\end{enumerate}

\begin{enumerate}
\def\labelenumi{\arabic{enumi}.}
\setcounter{enumi}{10}
\tightlist
\item
  设 A 为 Krull 整环，M、N 为 A 的分式域上有限秩 n 向量空间 V 中的两个有限生成格。对 A 中每个 \(c \neq 0\)（设它满足 \(cN \subset M\)），考虑行列式理想 \(\mathfrak{d}_{k}(M/cN)\)（习题 10）并记：
\end{enumerate}

\[
d _ {k} (\mathbf {N}, \mathbf {M}) = \operatorname{div} \left(\mathfrak {d} _ {k} (\mathbf {M} / c \mathbf {N})\right) - (n - k) \operatorname{div} (c).
\]

\begin{enumerate}
\def\labelenumi{(\alph{enumi})}
\item
  证明 \(d_{k}(\mathbf{N}, \mathbf{M})\) 不依赖于使 \(cN \subset M\) 的 c 的选取；\(d_{k}(\mathbf{N}, \mathbf{M})\) 称为 N 关于 M 的指标为 k 的行列式因子。
\item
  \(d_{k}(\mathbf{N},\mathbf{M})\geqslant d_{k + 1}(\mathbf{N},\mathbf{M})\) 且 \(d_{n}(\mathbf{N},\mathbf{M}) = 0\)。证明：若记
\end{enumerate}

\[
e _ {k} (\mathrm{N}, \mathrm{M}) = d _ {n - k} (\mathrm{N}, \mathrm{M}) - d _ {n - k + 1} (\mathrm{N}, \mathrm{M}),
\]

则

\[
e _ {k} (\mathbf {N}, \mathbf {M}) \leqslant e _ {k + 1} (\mathbf {N}, \mathbf {M})
\]

对 \(1 \leqslant k \leqslant n\) 成立；除子 \(e_k(\mathbf{N}, \mathbf{M})\) 称为 \(\mathbf{N}\) 关于 \(\mathbf{M}\) 的不变因子（归结为 \(\mathbf{A}\) 是主理想整环的情形）。

\begin{enumerate}
\def\labelenumi{(\alph{enumi})}
\setcounter{enumi}{2}
\item
  证明 \(d_0(\mathbf{N}, \mathbf{M}) = \chi(\mathbf{M}, \mathbf{N})\)（第 5 节）。
\item
  若 M 是 V 中的格而 N 是 V 的秩为 q \textless{} n 的子空间 W 中的格，则 N 关于 \(M \cap W\) 的不变因子称为 N 关于 M 的不变因子。说明《代数》第 VII 章 § 4 习题 8 至 10 与 14 至 16 可以怎样推广（在某些情形必要时假设 A 是整闭 Noether 整环）。
\end{enumerate}

\begin{enumerate}
\def\labelenumi{\arabic{enumi}.}
\setcounter{enumi}{11}
\tightlist
\item
  设 A 为整闭 Noether 整环，M、N 为两个同秩 r 的无扭有限生成 A-模，且 \(N \subset M\)；设 \(j: N \to M\) 为典范单射。
\end{enumerate}

\begin{enumerate}
\def\labelenumi{(\alph{enumi})}
\tightlist
\item
  证明对所有 \(k\)，\(\bigwedge^{k}j\colon\bigwedge^{k}\mathrm{N}\to\bigwedge^{k}\mathrm{M}\) 是伪单射的，\(\operatorname{Coker}\left(\bigwedge^{k}j\right)\) 是扭 A-模，且：
\end{enumerate}

\[
\chi \left(\operatorname{Coker} \left(\bigwedge^ {k} j\right)\right) = \binom {r - 1} {k - 1} \chi (\operatorname{Coker} j).
\]

\begin{enumerate}
\def\labelenumi{(\alph{enumi})}
\setcounter{enumi}{1}
\tightlist
\item
  利用 (a) 证明：若 M 是无扭有限生成 A-模，则 \(\bigwedge^{k}\mathbf{M}\) 的扭子模对所有 \(k\) 都是伪零的，且
\end{enumerate}

\[
c \left(\bigwedge^ {k} \mathbf {M}\right) = \binom {r - 1} {k - 1} c (\mathbf {M}),
\]

其中 \(\mathbf{M}\) 的秩为 \(r\)；特别地，\(c\left(\bigwedge^{r}\mathbf{M}\right) = c(\mathbf{M})\)。

\begin{enumerate}
\def\labelenumi{\arabic{enumi}.}
\setcounter{enumi}{12}
\tightlist
\item
  设 \(k\) 为域，\(A\) 为多项式环 \(k[X, Y]\)。
\end{enumerate}

\begin{enumerate}
\def\labelenumi{(\alph{enumi})}
\item
  设 \(m\) 为极大理想 \(AX + AY\)（在 \(A\) 中）；证明存在从 \(m\) 到 \(A\) 的伪同构，但不存在从 \(A\) 到 \(m\) 的伪同构。
\item
  设 \(\mathfrak{m}'\) 为极大理想 \(\mathbf{A}(\mathbf{X} - 1) + \mathbf{A}\mathbf{Y}\)（在 \(\mathbf{A}\) 中）；证明 \(c(\mathfrak{m}) = c(\mathfrak{m}') = 0\)，但既不存在从 \(\mathfrak{m}\) 到 \(\mathfrak{m}'\) 的伪同构，也不存在从 \(\mathfrak{m}'\) 到 \(\mathfrak{m}\) 的伪同构。
\item
  设 \(\mathfrak{p} = \mathbf{AX}\)，\(\mathfrak{q} = \mathbf{AY}\)，它们是高度为 1 的素理想。在 A-模 \(\mathbf{L} = \mathbf{A}^2\) 中考虑子模 \(\mathbf{M} = \mathfrak{pe}_1 \oplus \mathfrak{qe}_2\) 与 \(\mathbf{N} = \mathbf{Ae}_1 \oplus \mathfrak{pqe}_2\)（\(e_1, e_2\) 为 \(\mathbf{L}\) 的典范基的向量）。证明 \(\mathbf{M}\) 与 \(\mathbf{N}\) 同构，且存在从 \(\mathbf{L} / \mathbf{M}\) 到 \(\mathbf{L} / \mathbf{N}\) 的伪同构以及从 \(\mathbf{L} / \mathbf{N}\) 到 \(\mathbf{L} / \mathbf{M}\) 的伪同构，但不存在从 \(\mathbf{L}\) 到自身、把 \(\mathbf{M}\) 映到 \(\mathbf{N}\) 或把 \(\mathbf{N}\) 映到 \(\mathbf{M}\) 的伪同构（注意，借助命题 10，从 \(\mathbf{L}\) 到自身的伪同构必然是 \(\mathbf{L}\) 的自同构）。
\end{enumerate}

¶ 14. 设 A 为 Prüfer 整环（§ 2 习题 12），M、N 为 A 的分式域上有限秩 n 向量空间 V 中的两个有限生成格。对 A 中每个 \(c \neq 0\)（设它满足 \(cN \subset M\)），考虑行列式理想 \(\mathfrak{d}_{k}(M/cN)\)（习题 10）并记：

\[
\mathfrak {d} _ {k} (\mathbf {N}, \mathbf {M}) = c ^ {k - n} \mathfrak {d} _ {k} (\mathbf {M} / c \mathbf {N}).
\]

\begin{enumerate}
\def\labelenumi{(\alph{enumi})}
\item
  证明 \(\mathfrak{d}_k(\mathbf{N},\mathbf{M})\) 不依赖于 \(c\) 的选取，只要 \(c\mathbf{N}\subset \mathbf{M}\)；它们称为 \(\mathbf{N}\) 关于 \(\mathbf{M}\) 的行列式理想。
\item
  \(\mathfrak{d}_k(\mathbf{N},\mathbf{M})\subset \mathfrak{d}_{k + 1}(\mathbf{N},\mathbf{M})\) 且 \(\mathfrak{d}_n(\mathbf{N},\mathbf{M}) = \mathbf{A}\)。证明：若记 \(\mathfrak{e}_k(\mathbf{N},\mathbf{M}) = \mathfrak{d}_{n - k}(\mathbf{N},\mathbf{M})(\mathfrak{d}_{n - k + 1}(\mathbf{N},\mathbf{M}))^{-1}\)，则整理想 \(\mathfrak{e}_k(\mathbf{N},\mathbf{M})\) 满足 \(\mathfrak{e}_k(\mathbf{N},\mathbf{M})\supset \mathfrak{e}_{k + 1}(\mathbf{N},\mathbf{M})\)（对 \(1\leqslant k\leqslant n\)）；有限生成理想 \(\mathfrak{e}_k(\mathbf{N},\mathbf{M})\) 称为 \(\mathbf{N}\) 关于 \(\mathbf{M}\) 的不变因子（考虑 \(\mathrm{A_m}\) -模 \(\mathrm{M_m}\) 与 \(\mathrm{N_m}\)，其中 \(\mathfrak{m}\) 取遍 \(\mathrm{A}\) 的每个极大理想）。
\item
  若 M 是 V 的格而 N 是 V 的秩为 q \textless{} n 的子空间 W 中的格，则 N 关于 M ∩ W 的不变因子称为 N 关于 M 的不变因子。说明《代数》第 VII 章 § 4 习题 8 至 10 与 14 至 16 可以怎样推广。
\end{enumerate}

\begin{enumerate}
\def\labelenumi{\arabic{enumi}.}
\setcounter{enumi}{14}
\tightlist
\item
  设 A 为多项式环 Z{[}X, Y, T, U, V, W{]}，L 为自由 A-模 \(A^{4}\)，\((e_{i})_{1 \leq i \leq 4}\) 为其典范基，M 为 L 的由四个向量 \(Xe_{1}\)、\(Ye_{2}\)、\(Te_{3} + Ue_{4}\)、\(Ve_{3} + We_{4}\) 生成的子模。证明
\end{enumerate}

\[
\mathfrak {d} _ {1} (\mathrm{L} / \mathrm{M}) \mathfrak {d} _ {3} (\mathrm{L} / \mathrm{M}) \notin (\mathfrak {d} _ {2} (\mathrm{L} / \mathrm{M})) ^ {2}
\]

（与习题 14 (b) 比较）。

¶ 16. (a) 设 A 为 Prüfer 整环，M 为 A 的分式域 K 上有限秩 n 向量空间 V 中的有限生成格。证明存在 V 的基 \((e_{i})_{1 \leqslant i \leqslant n}\) 与 n 个分式理想 \(a_{i}\)（\(1 \leqslant i \leqslant n\)），使 M 等于诸 \(a_{i}e_{i}\) 的直和。（对 n 作归纳论证：若 \((u_{i})_{1\leq i\leq n}\) 是 V 的一组基，考虑有限生成分式理想 \(b_{1}\)（由 M 中元素的 \(u_{1}\) -坐标生成）；通过考虑 A-模 \(b_{1}^{-1}M\)，归结为 \(b_{1}=A\) 的情形。）

由此推出，若 W 是 V 的子空间，则 M ∩ W 是 M 的直和项。

\begin{enumerate}
\def\labelenumi{(\alph{enumi})}
\setcounter{enumi}{1}
\tightlist
\item
  设 M 是 \(A^{n}\) 的子模，等于诸 \(a_{i}e_{i}\) 的直和，其中诸 \(a_{i}\) 是 A 的有限生成（整）理想，且 \((e_{i})_{1\leqslant i\leqslant n}\) 是 \(A^{n}\) 的典范基。证明存在基 \((u_{i})_{1\leqslant i\leqslant n}\)（在 \(A^{n}\) 中）与 A 的理想 \(b_{i}\)，使 M 是诸 \(b_{i}u_{i}\) 的直和，且 \(b_{i}\) 整除 \(b_{i+1}\)（对 \(1\leqslant i\leqslant n-1\)）。（归结为 n=2 且 \(a_{1}+a_{2}=A\) 的情形。）
\end{enumerate}

\begin{enumerate}
\def\labelenumi{\arabic{enumi}.}
\setcounter{enumi}{16}
\tightlist
\item
  设 k 为域，A 为多项式环 \(k[X, Y]\)，L 为自由 A-模 \(A^{2}\)，\((e_{1}, e_{2})\) 为 L 的典范基。设 M 为 L 的由两个向量 \((X + 1)e_{1} + Ye_{2}, Ye_{1} + Xe_{2}\) 生成的子模。证明不存在从 L 到自身的伪同构，它限制在 M 上是从 M 到形如 \(au + bv\) 的子模 N 的伪同构（其中 \((u, v)\) 是 A-模 L 的基，a、b 是 A 的两个理想），或者限制在 N 上是从 N 到 M 的伪同构。（通过考虑行列式理想（习题 10）并注意 M 是自反的，可以只限于 N 也自反的情形，而此时必然有 a = A、b = AP，其中
\end{enumerate}

\[
\mathrm{P} (\mathrm{X}, \mathrm{Y}) = \mathrm{X} (\mathrm{X} + 1) - \mathrm{Y} ^ {2}
\]

是 A 的一个极端元；最后证明，对 L 的每组基 \((u, v)\)，M ∩ Au 既不能包含 Au，也不能包含于 APu。）

¶ 18. 设 A 为 Dedekind 整环，K 为其分式域，M、N 为 K 上有限秩 n 向量空间 V 中的两个格。设 \(\mathfrak{e}_{k} = \mathfrak{e}_{k}(N, M)\) 为 N 关于 M 的不变因子（习题 14）。证明存在 V 的基 \((u_{i})_{1 \leqslant i \leqslant n}\)，使 M 等于直和 \(\bigoplus_{i} a_{i} u_{i}\)，其中诸 \(a_{i}\) 是分式理想，而 N 等于直和 \(\bigoplus_{i} e_{i} a_{i} u_{i}\)。（利用主理想整环上有限生成模的理论（《代数》第 VII 章 § 4 第 2 节定理 1）以及单模群 \(\mathbf{SL}(n, \mathbf{A})\) 中的逼近定理（§ 2 第 4 节）。）

把这一结果推广到 A 是一族右有向的 Dedekind 整环子环之并的情形（例如 Dedekind 整环在其分式域的代数闭包中的整闭包）。

\begin{enumerate}
\def\labelenumi{\arabic{enumi}.}
\setcounter{enumi}{18}
\item
  设 A 为 Dedekind 整环，a、b 为 A 的分式理想。证明 A-模 A ⊕ ab 与 a ⊕ b 同构。
\item
  设 A 为 Dedekind 整环，P 为投射 A-模，N 为 P 的子模。证明 N 是投射的，且是同构于 A 的理想的模的直和。（可以只限于 P 是自由模的情形。然后照习题 16 那样进行，使用以《代数》第 VII 章 § 3 定理 1 为基础的超限归纳论证。）
\item
  设 P 为 Dedekind 整环 A 上的投射模。证明若 P 不是有限生成的，则它是自由的。（据习题 20，可把 P 写成无限直和 \(\bigoplus_{\lambda} (\mathfrak{b}_{\lambda} \oplus \mathfrak{c}_{\lambda})\)，其中，对每个 \(\lambda, \mathfrak{b}_{\lambda}\) 与 \(\mathfrak{c}_{\lambda}\) 都是 A 的理想。利用习题 19，有 \(\mathbf{P} = \mathbf{L} \oplus \mathbf{Q}\)，其中 \(\mathbf{L}\) 同构于 \(\mathbf{A}^{(I)}\)（I 为无限集），而 \(\mathbf{Q} = \bigoplus_{\alpha \in I} \mathfrak{a}_{\alpha}\)，其中诸 \(\mathfrak{a}_{\alpha}\) 是 A 的理想。然后应用《代数》第 II 章 § 2 习题 3。）
\end{enumerate}

¶ 22. 设 A 为局部环，m 为其极大理想，E 为有限生成 A-模。

\begin{enumerate}
\def\labelenumi{(\alph{enumi})}
\item
  若 \(r = \gamma(\mathrm{E})\) 是 \(\mathbf{E}\) 的一个生成系的最小基数，则 \(r\) 也是最小整数 \(h\)，使行列式理想 \(\mathfrak{d}_h(\mathbf{E})\) 等于 \(\mathbf{A}\)，同时也是最小整数 \(k\)，使 \(\bigwedge^{k+1} \mathbf{E} = 0\)（注意 \(r\) 是 \(\mathbf{E}/\mathfrak{m}\mathbf{E}\) 在 \(\mathbf{A}/\mathfrak{m}\) 上的秩）。
\item
  若记 \(e(E) = \mathfrak{d}_{r-1}(E)\)，证明 \(\bigwedge^{r} E\) 同构于 \(A / e(E)\)。理想 \(a\)（在 \(A\) 中）包含 \(e(E)\) 的充分必要条件是 \(E / aE\) 是自由 \((A / a)\) -模（归结为 \(a = 0\) 的情形）。
\item
  若 \(\mathfrak{e}(\mathbf{E})\) 是主理想 \(A\alpha\)，证明 E 包含一个同构于 \(A/A\alpha\) 的直和项（把 E 写成 \(A^{r}\) 对某个子模 R 的商，并注意 R 中有一个形如 \(\alpha y\) 的元素，其中 y 是 \(A^{r}\) 的某组基中的元素）。
\item
  设 \(\mathfrak{b}_h^{\prime}(\mathbf{E})\) 为 \(\bigwedge^{h+1}\mathbf{E}\) 的零化子（对 \(0 \leqslant h \leqslant r - 1\)）。证明下列条件等价：
\end{enumerate}

\((\alpha)\) 诸理想 \(\mathfrak{d}_h(\mathbf{E})\)（\(0 \leqslant h \leqslant r - 1\)）是主理想。

(β) 诸理想 \(\mathfrak{d}_h^\prime (\mathrm{E})\) \((0\leqslant h\leqslant r - 1)\) 是主理想。

\((\gamma)\) E 是 r 个同构于 A/A \(\lambda_{h}\) 的模的直和，其中 \(\lambda_{h+1}\) 整除 \(\lambda_{h}\)（对 \(0 \leqslant h \leqslant r - 1\)）。

而且，若这些条件成立，则 \(\mathfrak{d}_h'(\mathrm{E}) = \mathrm{A}\lambda_h\)（对 \(0 \leqslant h \leqslant r - 1\)），且 \(\mathfrak{d}_h(\mathrm{E}) = \lambda_h\lambda_{h+1} \ldots \lambda_{r-1}\mathrm{A}\)（对 \(0 \leqslant h \leqslant r - 1\)）。（对 \(r\) 作归纳，并利用 (c)。）

\begin{enumerate}
\def\labelenumi{(\alph{enumi})}
\setcounter{enumi}{4}
\tightlist
\item
  再设 A 是整环。证明：若 p 是 A 的素理想，使 E/pE 是自由 (A/p)-模且 \(E_{p}\) 是自由 \(A_{p}\) -模，则 E 是自由 A-模（利用 (b) 证明 \(\gamma(E_{p}) = \gamma(E)\) 且 \(e(E_{p}) = e(E)\)）。
\end{enumerate}

\begin{enumerate}
\def\labelenumi{\arabic{enumi}.}
\setcounter{enumi}{22}
\item
  设 A 为环，E 为有限生成 A-模，r 为使 \(\bigwedge^{k+1}E=0\) 的最小整数 k，\(e(E)\) 为 \(\bigwedge^{r}E\) 的零化子。证明对 A 中每个理想 \(a\supset e(E)\)，\((A/a)\) -模 E/aE 是平坦的（归结为 A 是局部环的情形并利用习题 22 (b)）。由此推出：若 r 是 A 的 Jacobson 根，且对 A 的每个极大理想 m，向量 \((A/m)\) -空间 E/mE 的秩都相同，则 E/rE 是平坦 \((A/r)\) -模。
\item
  设 A 为整闭 Noether 整环。格 M（关于 A 的）是自反的，其充分必要条件是它满足下列条件：对 A 中元素的每个有序对 \((a, b)\)，若 \(a \neq 0\) 且 A/aA 上比率为 b 的位似是单射，则 M/aM 上比率为 b 的位似也是单射。（为看条件必要，考虑 M 中满足 \(ax + by = 0\) 的两个元素 x、y，并证明对所有 \(p \in P(A)\) 有 \(y \in aM_{p}\)。为证条件充分，利用定理 2 的判别法 (c)，并注意（利用 § 1 第 4 节命题 8）假设可以写成
\end{enumerate}

\[
\operatorname{Ass} (\mathbf {M} / a \mathbf {M}) = \operatorname{Ass} (\mathbf {A} / a \mathbf {A})
\]

对 \(a \neq 0\)（\(\mathbf{A}\) 中）成立，且对每个模 \(\mathbf{E}\)（满足 \(\mathbf{M} \subset \mathbf{E} \subset \mathbf{V}\)），存在 \(a \neq 0\)（在 \(\mathbf{A}\) 中）使 \(a\mathbf{M} \subset a\mathbf{E} \subset \mathbf{M}\)。

\begin{enumerate}
\def\labelenumi{\arabic{enumi}.}
\setcounter{enumi}{24}
\tightlist
\item
  设 \(A = \bigoplus_{n \geq 0} A_n\) 为正次数的 Noether 分次环。证明：若 A 是因子分解整环，则 \(A_0\) 是因子分解整环，且诸 \(A_n\) 是自反的。（为证 \(A_0\) 是因子分解整环，利用 § 3 第 2 节定理 1 的判别法 (c)；为证诸 \(A_n\) 自反，利用习题 24。）
\end{enumerate}

¶ 26. 设 A 为 Noether 环，M 为有限生成 A-模。对称代数 S(M) 是因子分解整环的充分必要条件是：A 是因子分解整环且诸 \(S^{n}(M)\) 是自反 A-模。（为证条件充分，首先注意，若 T = A - \{0\}，则 S(M) 等同于 \(T^{-1}S(M)\) 的一个子环；然后证明 A 的每个极端元 p 在 S(M) 中也是极端元，这可归结为证明：若 p 整除 S(M) 中两个齐次元素之积 xy，则它整除二者之一；最后应用 § 3 第 4 节命题 3。）

